% !PRINTED 12
\weberpage{12}
%<<CONT
than, equal to, or greater than $\mu$. If $e'$ is equal to $\mu$, then
rational ratios lie both between $e$ and $e'$ and between $e'$ and
$\epsilon$. If, however, $e'$ is smaller than $\mu$, then $\mu$ lies
between $e'$ and $\epsilon$, and if $e'$ is greater than $\mu$, then
$\mu$ lies between $e$ and $e'$; hence $e'$ cannot at the same time be
equal to $e$ and equal to $\epsilon$.

In case 2) the ratios $e,\epsilon$ are called \gesperrt{unequal}. Thus
either $\alpha)$ $e<\mu<\epsilon$ or $\beta)$
$e>\mu'>\epsilon$ may hold, and these two cases exclude one another,
because from $e<\mu<\epsilon$, $\epsilon<\mu'$ it follows that
$\mu<\mu'$, and consequently $e<\mu'$ results.

The relation of magnitude $2\alpha)$ or $2\beta)$ also remains
unchanged when $e$ or $\epsilon$ is replaced by an element equal to it.
For if $\mu<\epsilon$ and $\mu\gtrless\epsilon'$, then a rational ratio
lies between $\epsilon$ and $\epsilon'$, and $\epsilon,\epsilon'$ are
not equal.

In case $2\alpha)$, $e$ is called \gesperrt{smaller than} $\epsilon$;
in case $2\beta)$, $e$ is called \gesperrt{greater than} $\epsilon$.
Between two unequal ratios one can insert an arbitrary number of
rational ratios.

If now we collect all ratios equal to one another, we obtain a generic
concept which we designate as a \gesperrt{number} in the general sense
of the word. A number is therefore a name or sign for a certain
manifold whose elements are precisely the ratios equal to one of them.
\setcounter{footnote}{0}\footnote{The natural numbers too can be traced back in a simple and
consistent way to the generic concept.} Under this concept of number
are included the rational ratios, and hence also the natural numbers as
the ratios $m:1$; they form the \gesperrt{rational numbers}. Numbers
which do not arise from rational ratios are called \gesperrt{irrational
numbers}. According to what has been developed so far, the numbers form
an \gesperrt{ordered set}, and their order can be fixed if, for each
number, any one of the ratios contained under it is chosen as
representative.

\begin{center}
\rule{0.18\linewidth}{0.4pt}
\end{center}

Of two ratios with the same denominator and unequal numerators, that
one is greater whose numerator is greater; and
%CONT>>
